% Компонентная диаграмма гибридной интеллектуальной системы.
\begin{tikzpicture}[
  font=\footnotesize,
  box/.style={draw, rounded corners=2pt, align=center, inner sep=3pt, minimum height=9mm},
  ext/.style={box, dashed},
  arr/.style={-{Stealth[length=2mm]}, semithick},
  biarr/.style={{Stealth[length=2mm]}-{Stealth[length=2mm]}, semithick},
  lbl/.style={font=\scriptsize, align=center, inner sep=1pt},
]

% --- клиент виртуальной среды ---
\node[box, minimum width=2.3cm] (chat) at (0, 0) {Игровой чат};
\node[box, minimum width=2.3cm] (overlay) at (0, -1.7) {Оверлей-экран};
\node[box, minimum width=2.3cm] (state) at (0, -3.4) {Сбор состояния\\игрока};
\node[draw, thick, inner sep=2.5mm, fit=(chat)(overlay)(state),
      label={[font=\footnotesize]above:Клиент среды}] (client) {};

% --- сервер ---
\node[box, minimum width=5.2cm] (orch) at (7.0, 0) {Оркестратор\\конвейер обработки реплики};
\node[box, minimum width=2.1cm, minimum height=14mm] (resolver) at (5.45, -2.7) {Резолвер};
\node[box, minimum width=2.1cm, minimum height=14mm] (solver) at (8.55, -2.7)
  {Решатель\\задач};
\node[box, minimum width=5.2cm] (kb) at (7.0, -5.3) {База знаний\\онтология RDF/OWL};
\node[draw, thick, inner sep=2.5mm, fit=(orch)(resolver)(solver)(kb),
      label={[font=\footnotesize]above:Сервер}] (server) {};

% --- внешние системы ---
\node[ext, minimum width=2.4cm] (llm) at (13.3, 0) {Провайдер\\языковой модели};
\node[box, minimum width=2.4cm] (importer) at (13.3, -5.3) {Импортёр};
\node[ext, minimum width=2.4cm] (data) at (13.3, -7.2) {Данные среды\\и локализация};

% --- клиент и сервер ---
\draw[arr] ([yshift=3mm]orch.west -| client.east) -- node[lbl, above] {реплика,\\состояние}
  ([yshift=3mm]orch.west);
\draw[arr] ([yshift=-3mm]orch.west) -- node[lbl, below] {ответ, план,\\команда}
  ([yshift=-3mm]orch.west -| client.east);

% --- оркестратор и модель ---
\draw[arr] ([yshift=3mm]orch.east) -- node[lbl, above] {реплика}
  ([yshift=3mm]llm.west);
\draw[arr] ([yshift=-3mm]llm.west) -- node[lbl, below] {намерение}
  ([yshift=-3mm]orch.east);

% --- внутри сервера: подписи собраны в промежутке между колонками ---
\draw[biarr] (orch.south -| resolver.north) -- node[lbl, right] {название /\\сущность}
  (resolver.north);
\draw[biarr] (orch.south -| solver.north) -- node[lbl, left] {действие /\\вердикт,\\план}
  (solver.north);
\draw[biarr] (resolver.south) -- node[lbl, right] {имена} (resolver.south |- kb.north);
\draw[biarr] (solver.south) -- node[lbl, left] {рецепты,\\источники,\\запреты}
  (solver.south |- kb.north);

% --- сборка базы знаний ---
\draw[arr] (data.north) -- (importer.south);
\draw[arr] (importer.west) -- node[lbl, above, pos=0.35] {триплеты} (kb.east);

\end{tikzpicture}
